\documentclass[11pt,a4paper]{article}

% ==== Fonts & CJK ====
\usepackage{xeCJK}
\setCJKmainfont{Songti SC}
\setCJKsansfont{PingFang SC}
\setCJKmonofont{Menlo}
\setmainfont{Times New Roman}
\setmonofont{Menlo}

\usepackage[margin=2.2cm]{geometry}
\usepackage{amsmath, amssymb, amsthm, mathtools}
\usepackage{bm}
\usepackage{booktabs, tabularx, array, multirow}
\usepackage{graphicx}
\usepackage{xcolor}
\usepackage{hyperref}
\usepackage{enumitem}
\usepackage{caption}
\usepackage{fancyhdr}
\usepackage{tcolorbox}

\hypersetup{
  colorlinks=true,
  linkcolor=blue!60!black,
  urlcolor=blue!60!black,
  citecolor=blue!60!black,
  pdftitle={RS+BCH 级联码课题——需求澄清后的技术方案},
  pdfauthor={陈诗阳}
}

\pagestyle{fancy}
\fancyhf{}
\rhead{\footnotesize RS+BCH 级联码课题}
\lhead{\footnotesize 需求澄清后的技术方案}
\cfoot{\thepage}

\renewcommand{\baselinestretch}{1.2}

\definecolor{accent}{RGB}{30,80,160}
\definecolor{accent2}{RGB}{160,40,40}
\definecolor{accent3}{RGB}{40,120,60}

\newtcolorbox{keybox}[1][]{colback=blue!5, colframe=accent, boxrule=0.4mm,
  arc=1mm, left=2mm, right=2mm, top=1mm, bottom=1mm, #1}
\newtcolorbox{riskbox}[1][]{colback=orange!5, colframe=orange!70!black, boxrule=0.4mm,
  arc=1mm, left=2mm, right=2mm, top=1mm, bottom=1mm, #1}
\newtcolorbox{actionbox}[1][]{colback=green!5, colframe=accent3, boxrule=0.4mm,
  arc=1mm, left=2mm, right=2mm, top=1mm, bottom=1mm, #1}
\newtcolorbox{quotebox}[1][]{colback=gray!8, colframe=gray!50, boxrule=0.3mm,
  arc=1mm, left=2mm, right=2mm, top=1mm, bottom=1mm, #1}

\begin{document}

% ============ Cover ============
\begin{center}
{\Large\bfseries RS+BCH 级联码课题}\\[4pt]
{\large 需求澄清后的技术方案与执行计划}\\[16pt]
{\normalsize 陈诗阳\\
基于 IEEE TIT 2025 论文《Efficient Ordered Statistics Decoding of BCH Codes Without Gaussian Elimination》\\
\vspace{2pt}
\texttt{https://github.com/a-yeyang/efficient-osd-bch-no-ge}}\\[10pt]
\rule{0.9\textwidth}{0.4pt}\\[6pt]
{\footnotesize 2026-07-23}\\
\rule{0.9\textwidth}{0.4pt}
\end{center}

\vspace{6pt}
\begin{keybox}[title=\textbf{文档目的}]
根据老师对之前 6 个对齐问题 (Q1--Q6) 的答复 (R1--R6)，
本文档整理出：
\begin{enumerate}[leftmargin=1.5em, itemsep=1pt, topsep=2pt]
\item 逐条回答的\textbf{技术解读}和对整体方案的\textbf{影响}
\item 因回答衍生出的\textbf{三个新技术追问}，请老师再指点
\item 更新后的\textbf{工作量估算和分阶段执行计划}
\end{enumerate}
\end{keybox}

% ============ Section 1 ============
\section{老师答复的逐条解读}

\subsection{R1 —— 内码单独 / 内外码都用 Lagrange 两种方案并列}

\begin{quotebox}
\textbf{R1 原文}：请尝试内码单独、内外码两者都用 Lagrange 插值的方案。
\end{quotebox}

\textbf{解读}：R1 确认了``Lagrange 共享''是课题的核心创新点，并明确要求做\textbf{消融对比 (ablation)}
来量化共享带来的收益：

\begin{table}[!ht]
\centering
\begin{tabularx}{\textwidth}{lXXX}
\toprule
\textbf{方案} & \textbf{内码 Lagrange} & \textbf{外码 Lagrange} & \textbf{共享代数结构} \\
\midrule
Baseline & --- & 否 (BM) & 无 \\
方案 A & 是 (LLOSD) & 否 (LCC-BR 传统实现) & 部分 \\
方案 B & 是 (LLOSD) & 是 (LCC-BR + Lagrange 共享) & 全共享 \\
\bottomrule
\end{tabularx}
\caption{R1 引出的三级消融实验}
\end{table}

\textbf{对整体方案的影响}：
\begin{itemize}[leftmargin=1.5em, itemsep=1pt]
\item 需要\textbf{多做一个外码译码器}：Lagrange 版本的 RS 软判 (与内码 LLOSD 共享 $\alpha^j$ 表、
denominator products、Lagrange 基函数 $T_j(\alpha^i)$)。
\item 汇报时对比：Baseline / 方案 A / 方案 B 三条时延曲线，从共享粒度上量化收益。
\item 预期效果：方案 B 相对方案 A 能省去 $O(n\cdot k')$ 次 $\mathbb{F}_{2^m}$ 乘法
（Lagrange 基函数表的重建成本）。
\end{itemize}

\subsection{R2 —— 10\% 时延 KPI 对两个基线都要满足}

\begin{quotebox}
\textbf{R2 原文}：10\% 时延是分别 RS 硬判决和软判决的时延基线。针对时钟周期和 $\mathbb{F}_{2^m}$ ops。
\end{quotebox}

\textbf{解读}：这条明确了 KPI 的具体形式：

\begin{keybox}[title=\textbf{硬性 KPI}]
\begin{itemize}[leftmargin=1.5em, itemsep=1pt, topsep=2pt]
\item $T_{\text{RS+BCH}} \leq 1.10 \times T_{\text{RS-BM}}$ \hfill (Case a：硬判基线)
\item $T_{\text{RS+BCH}} \leq 1.10 \times T_{\text{RS-LCC-BR}}$ \hfill (Case b：软判基线)
\end{itemize}
两个 case 都需要满足，且两种测量口径 (时钟周期 \& $\mathbb{F}_{2^m}$ ops) 都要报告。
\end{keybox}

\textbf{初步分析（基于 IEEE TIT 2025 的 Table IV 及 PLCC 论文的实测）}：

\begin{table}[!ht]
\centering
\begin{tabularx}{\textwidth}{Xllll}
\toprule
\textbf{译码器} @ 5 dB & \textbf{$\mathbb{F}_{2^m}$ ops} & \textbf{相对 BM 倍数} & \textbf{是否满足 Case a} & \textbf{是否满足 Case b} \\
\midrule
RS-BM (n=255)              & $\sim 3\times 10^3$   & 1.0$\times$  & --- & --- \\
RS-LCC-BR (n=255)          & $\sim 3.75\times 10^5$ & 125$\times$  & --- & --- \\
BCH-LLOSD(2) (n=255)       & $\sim 1.9\times 10^4$  & 6.3$\times$  & \textcolor{red}{$\times$ ($>$10\%)}  & \textcolor{accent3}{$\checkmark$ ($\ll$10\%)} \\
BCH-HSD(1,8) (n=255)       & $\sim 9.9\times 10^3$  & 3.3$\times$  & \textcolor{red}{$\times$ (需极致优化)} & \textcolor{accent3}{$\checkmark$ ($\ll$10\%)} \\
\bottomrule
\end{tabularx}
\caption{两种 KPI 基线下 BCH 内码时延占比初步估算 (n=255)}
\end{table}

\begin{riskbox}[title=\textbf{Case a 的可行性分析}]
\textbf{结论 A}：Case b (相对 RS 软判基线) 满足 10\% 预算相对容易，因为 LCC-BR 本身开销
就是 BM 的 $\sim$100$\times$，BCH 加进去 (相当于 BM 的 $\sim$3--6$\times$) 只占 3--6\%，
留有充足 buffer。

\textbf{结论 B}：Case a (相对 RS 硬判基线) 从 F$_{2^m}$ ops 层面\textbf{几乎不可能满足}——
LLOSD 的运算量本质上是 BM 的 3--6$\times$，10\% 的预算需要压缩到 BM 的 0.1$\times$，
差 30$\times$ 以上。可能的应对：

\begin{itemize}[leftmargin=1em, itemsep=1pt, topsep=2pt]
\item 用极致的\textbf{硬件并行度} (P=64$\times$ 并行 F$_{2^m}$ 乘法器) 才有希望
\item 或者转换视角：在\textbf{时钟周期 (关键路径深度)}下测量而非总 ops 数
\item 或者只在\textbf{高 SNR (6 dB+)} 区间达到 (SNR 越高 LLOSD 越快，Table I 显示 6 dB 时可到 BM 的 8\%)
\end{itemize}
\end{riskbox}

\subsection{R3 —— n=128 和 n=255 并列做}

\begin{quotebox}
\textbf{R3 原文}：并列做。可以。
\end{quotebox}

\textbf{解读}：两个码长都要做，工作量翻倍。但真正的挑战在于两个 config 底层参数完全不同：

\begin{table}[!ht]
\centering
\begin{tabularx}{\textwidth}{lXX}
\toprule
& \textbf{n=128 方案} & \textbf{n=255 方案} \\
\midrule
Galois 域 & GF($2^7$) = GF(128) & GF($2^8$) = GF(256) \\
primitive polynomial & 例如 0x83 & 例如 0x11D \\
BCH 类型 & 需 \textbf{extended BCH} (127+1) 或 shortened BCH & primitive BCH (255 = $2^8 - 1$) \\
Lagrange 表大小 & $128 \times 128$ & $256 \times 256$ \\
存储 & $\sim$16 KB & $\sim$64 KB \\
\bottomrule
\end{tabularx}
\end{table}

\begin{riskbox}[title=\textbf{n=128 的技术风险 (需要老师定夺)}]
论文 LLOSD 的 Theorem 2 只覆盖 \textbf{primitive narrow-sense BCH}
(码长 $n = 2^m - 1$)。因此 $n=128$ 需要用：

\begin{itemize}[leftmargin=1em, itemsep=1pt, topsep=2pt]
\item \textbf{Extended BCH (加一位全奇偶校验)}：需要\textbf{推广 Theorem 2}
到扩展码——这是一个小 novelty (需要小心处理 syndrome 计算和过滤条件)
\item \textbf{Shortened BCH (删除信息位)}：更简单，但码率下降需要重新调整参数
\end{itemize}

无论哪种做法，都比 $n=255$ 多做一份专门工作。
\end{riskbox}

\subsection{R4 —— 外码软判用 LCC-BR}

\begin{quotebox}
\textbf{R4 原文}：用后者 (LCC-BR)。
\end{quotebox}

\textbf{解读}：LCC-BR 是最合适的外码软判选择，理由：

\begin{itemize}[leftmargin=1.5em, itemsep=1pt]
\item \textbf{与 LLOSD 同源}：论文 §V 的 HSD 算法就是 LLOSD + LCC-BR 集成，
两者共享 basis reduction 结构 —— \textbf{Lagrange 共享的技术路径已经在论文里被打通}。
\item \textbf{复杂度低、硬件友好}：Xing-Chen-Bossert 2020 论文中 LCC-BR 的 module basis
reduction 比 K\"otter–Vardy 的 factorization approach 少约 1/3 的 $\mathbb{F}_{2^m}$ 运算。
\item \textbf{与老师的时延 KPI 契合}：LCC-BR 的低时延特性正是我们要利用的。
\end{itemize}

\textbf{工作量}：LCC-BR 的核心是 Mulders--Storjohann basis reduction + partial root-finding。
Python 复现里当前是简化实现 (用 BM 替代)，Matlab 版\textbf{需要完整实现}，
预计工作量 1--1.5 周。

\subsection{R5 —— 所有译码器自己实现}

\begin{quotebox}
\textbf{R5 原文}：所有译码器自己实现。
\end{quotebox}

\textbf{解读}：不允许使用 Matlab Communications Toolbox 的
\texttt{comm.RSDecoder} / \texttt{comm.BCHDecoder} 等任何内置译码器。

\textbf{好消息}：Python 复现代码中所有的算法逻辑都可以 port，工作量集中在语言转换：

\begin{table}[!ht]
\centering
\begin{tabularx}{\textwidth}{lXX}
\toprule
Python 模块 & 对应 Matlab 模块 & 转换成本 \\
\midrule
\texttt{src/gf.py} & GF class + EXP/LOG 查表 & 低 (纯查表) \\
\texttt{src/bch.py} & BCH encoder + BM decoder & 低 (向量化直接对应) \\
\texttt{src/osd.py} & OSD baseline (GE + TEP enum) & 中 (Matlab GE 有内置) \\
\texttt{llosd.py} + \texttt{llosd\_jit.py} & LLOSD + 内循环 & 高 (Numba JIT 需用 MEX 或 vectorization 替代) \\
\bottomrule
\end{tabularx}
\end{table}

\textbf{预计整体 port 工作量}：2--3 周 (含单元测试与 Python 版结果对齐)。

\subsection{R6 —— 三个参数方案都需要}

\begin{quotebox}
\textbf{R6 原文}：三个方案都需要。
\end{quotebox}

\textbf{解读}：仿真矩阵扩展如下：

\begin{table}[!ht]
\centering
\begin{tabularx}{\textwidth}{lXll}
\toprule
方案 & (n, RS, BCH) & 码率 & 域 \\
\midrule
n=255 方案 A & RS(255, 239) + BCH(255, 239) & 0.879 & GF($2^8$) \\
n=255 方案 B & RS(255, 235) + BCH(255, 247) & 0.891 & GF($2^8$) \\
n=128 方案   & RS(127, 119) + BCH(127, 120) & 0.885 & GF($2^7$) \\
\bottomrule
\end{tabularx}
\end{table}

\textbf{仿真矩阵规模}：3 参数 $\times$ 2 基线 (BM / LCC-BR) $\times$ 2 方案 (A/B) $\times$
2 测量口径 (cycles / ops) $\times$ 15 SNR points = \textbf{约 360 个数据点}。
每个数据点约需 $10^4$ 帧收敛 $\to$ \textbf{$\sim 4\times 10^6$ 帧总仿真量}。

\textbf{应对}：使用 Matlab \texttt{parfor} 并行 + 参数化 config file 统一驱动。

% ============ Section 2 ============
\section{基于 R1--R6 的三个新技术追问}

在动手写代码前，还有三个技术细节希望老师再指点，因为它们会决定后续几周的具体路线：

\begin{actionbox}[title=\textbf{Q7. Clock cycle 测量的 hardware timing model}]
纯软件仿真无法直接得到 clock cycles。我建议定义一个\textbf{参数化的抽象硬件模型}：

\begin{itemize}[leftmargin=1em, itemsep=1pt, topsep=2pt]
\item 每个 $\mathbb{F}_{2^m}$ 乘法 = 1 cycle
\item 允许 $P$ 路并行 (硬件并行度)
\item 报告 $P \in \{1, 4, 16, 64\}$ 四种硬件档次下的关键路径 depth
\end{itemize}

请问这个 model 可以接受吗？还是老师有指定的 timing model？
\end{actionbox}

\begin{actionbox}[title=\textbf{Q8. Case (a) 的可行性策略}]
如 §1.2 分析，Case (a) 相对 RS 硬判决 BM 的 10\% 时延预算在 $\mathbb{F}_{2^m}$ ops 层面
几乎不可能满足 (需要将 LLOSD 时延压到 BM 的 0.1$\times$，差 30$\times$ 以上)。

请问在这种情况下：
\begin{itemize}[leftmargin=1em, itemsep=1pt, topsep=2pt]
\item [(a)] 是否可以只在\textbf{时钟周期 (关键路径)} 视角下 (即高并行度硬件模型) 满足 10\%？
\item [(b)] 是否可以只在\textbf{中高 SNR ($\geq 5.5$ dB)} 区间满足？
\item [(c)] 还是必须两个基线的 10\% 都要在所有 SNR、所有测量口径下无条件满足？
\end{itemize}

这个方向会决定我要不要花大量精力在极致优化上 (例如 Lagrange 查表化、TEP 循环展开等)。
\end{actionbox}

\begin{actionbox}[title=\textbf{Q9. n=128 用 Extended BCH 还是 Shortened BCH}]
$n=128$ 不是 primitive BCH 长度 ($2^m - 1$)，需要通过 extended 或 shortened 构造。
两种方式都会对 LLOSD 的 Theorem 2 需要小幅推广：

\begin{itemize}[leftmargin=1em, itemsep=1pt, topsep=2pt]
\item Extended BCH：$n = 128 = 127 + 1$，加一位全奇偶校验。
Theorem 2 的过滤条件需要\textbf{加一个额外的 syndrome 检查}。
\item Shortened BCH：例如 BCH(127, 119) 用作 (128, ?) 效果不佳；或从 (255, ?) shortened 到 128。
\end{itemize}

请问老师倾向哪种？\textbf{Extended BCH 是更规范但需要小 novelty；shortened 更简单但码率折损}。
\end{actionbox}

% ============ Section 3 ============
\section{更新后的工作量估算与执行计划}

\begin{table}[!ht]
\centering
\begin{tabularx}{\textwidth}{llX}
\toprule
\textbf{阶段} & \textbf{时长} & \textbf{交付物} \\
\midrule
Phase 1 & Week 1--2 &
Matlab 项目框架 + GF($2^7$)/GF($2^8$) 有限域算术 + BCH 编码
+ BM 硬判 + PAM4 modem + AWGN 信道模型。BER-SNR 基线曲线与文献对齐 \\
\midrule
Phase 2 & Week 3--4 &
LLOSD 从 Python 复现代码 port 到 Matlab (含 SLLOSD 和 ML 提前终止)。
BCH 单码仿真结果与原论文 Fig 3/4 对齐 \\
\midrule
Phase 3 & Week 5--6 &
RS-LCC-BR 完整实现 (Mulders-Storjohann basis reduction + partial root-finding)。
外码 RS 单码性能验证 \\
\midrule
Phase 4 & Week 7--8 &
级联链路搭建 + \textbf{Lagrange 共享层} (方案 A 和方案 B)。
n=255 方案 A 完整仿真 (最小闭环) \\
\midrule
Phase 5 & Week 9--10 &
横向扩展：n=255 方案 B + n=128 方案。
Extended BCH 的 Theorem 2 推广 (若采用) \\
\midrule
Phase 6 & Week 11--12 &
两种测量口径的时延统计、10\% KPI 验证、
最终报告 + Matlab 代码整理与交付 \\
\bottomrule
\end{tabularx}
\end{table}

\textbf{总工作量估计：10--12 周}。相比初版的 6--8 周，主要增量来自：

\begin{itemize}[leftmargin=1.5em, itemsep=1pt]
\item R6 的\textbf{三个参数方案}$\to$ 仿真组合数 $\times 3$
\item R3 的 n=128 需要\textbf{extended/shortened BCH 处理}
\item R1 的\textbf{方案 A / 方案 B 消融对比}
\item R2 的\textbf{两个基线} (BM \& LCC-BR) $\times$ 两种测量口径 (cycles \& ops)
\end{itemize}

\vspace{4pt}
\begin{keybox}[title=\textbf{建议的增量执行策略}]
先做\textbf{最小闭环 (n=255 方案 A + LCC-BR baseline + F$_{2^m}$ ops)}，
跑通后请老师 review 中间结果和路线正确性。确认后再横向扩展到其他 config。
避免同时展开所有维度导致后期发现底层 bug 需要全面 rework。

\vspace{4pt}
\hrule
\vspace{3pt}
{\footnotesize
\textbf{参考}: L.\ Yang et al., ``Efficient OSD of BCH Codes Without GE,''
{\it IEEE Trans.\ Inf.\ Theory}, vol.\ 71, no.\ 11, pp.\ 8294--8311, Nov 2025.
Python 复现: \url{https://github.com/a-yeyang/efficient-osd-bch-no-ge}.
文档: 2026-07-23 · 陈诗阳
}
\end{keybox}

\end{document}
